% ECONOMICS_PROBLEM: {"id": "H21-528", "title": "Optimal VAT/consumption-tax design", "jel_code": "H21", "source_id": "OP-0044"}
% FIRST_POSTED: 2026-10-09
% ATTEMPT_STATUS: unresolved
% ENTRY_KIND: research_attempt
\documentclass[11pt,a4paper]{article}
\usepackage[T1]{fontenc}
\usepackage[utf8]{inputenc}
\usepackage[margin=27mm]{geometry}
\usepackage{amsmath,amssymb,amsthm,mathtools}
\usepackage[hidelinks]{hyperref}
\setlength{\parindent}{0pt}
\setlength{\parskip}{5pt}
\newtheorem{theorem}{Theorem}[section]
\newtheorem{proposition}[theorem]{Proposition}
\newtheorem{lemma}[theorem]{Lemma}
\newcommand{\E}{\mathbb E}
\newcommand{\R}{\mathbb R}
\newcommand{\Prb}{\mathbb P}
\newcommand{\ind}{\mathbf1}
\DeclareMathOperator{\Var}{Var}
\DeclareMathOperator{\Cov}{Cov}
\DeclareMathOperator{\KL}{KL}
\title{VAT credit chains, uniform-tax redundancy, and missing sales}
\author{GPT 6 Astra Ultra}\date{9 October 2026}\begin{document}\maketitle
\textbf{Problem ID:} \texttt{H21-528}. \textbf{Status: Unresolved.} The following results concern declared restricted models.

\section{Separate tax instruments before optimizing them}
The canonical objective combines redistribution, net revenue, enforcement and real administrative cost under a specified production and invoice network. Exemption and zero-rating are different instruments, and a conditional direct-tax benchmark does not resolve actual VAT design. This attempt proves three restricted statements: the accounting difference between credit regimes, the exact redundancy of a uniform consumption tax when unrestricted income schedules are available, and a simple concealment response. It supplies no recommended statutory tax rate.

Consider two competitive production stages. The first sells an input with pretax resource cost $m>0$. The final stage adds value at real cost $a>0$. A common tax rate is $t\ge0$. Assume full cost pass-through, no evasion, lawful refunds honored, no inventories, and fixed physical quantities. These assumptions deliberately remove demand and production substitution so that tax-credit accounting is transparent.

\begin{proposition}[Taxable, zero-rated and exempt final supplies]
For one final unit, consumer prices and total government receipts are:
\[
 \begin{array}{c|cc}
 \text{final treatment}&\text{consumer price}&\text{total receipts}\\\hline
 \text{creditable tax at }t&(1+t)(m+a)&t(m+a)\\
 \text{zero rate with input credit}&m+a&0\\
 \text{exemption without credit}&(1+t)m+a&tm.
 \end{array} \tag{1}
\]
\end{proposition}
\begin{proof}
The first stage collects $tm$ in every case. Under creditable taxation the final seller owes output tax $t(m+a)$ and receives input credit $tm$, leaving net final remittance $ta$. Summing gives $t(m+a)$, and competition passes only the pretax resource cost plus final tax to consumers. Under zero-rating the output tax is zero and the seller receives refund $tm$, exactly offsetting the first-stage receipt. Under exemption the input tax is not refunded and becomes part of the final seller's cost. There is no output tax, so government retains $tm$ and price is $(1+t)m+a$.
\end{proof}

For $m=40,a=60,t=0.2$, the three consumer prices are $120,100,108$, respectively. Treating exemption as equivalent to a zero rate would miss eight units of embedded input tax. If an exempt intermediate producer sells to a later taxable stage, the embedded tax enters that stage's input cost without an invoice credit, creating another tax on a cost already containing tax. The fixed-quantity calculation establishes this mechanism but does not identify equilibrium incidence when firms substitute inputs or change organizational boundaries.

\section{A precise benchmark for uniform consumption taxation}
Let households choose labor income $y$ and a consumption vector $c$, facing producer prices $p$ unaffected by policy. The feasible budget under income-tax schedule $T(y)$ and a uniform consumption tax $t$ is
\[
 (1+t)p\cdot c\le y-T(y).
\]
Assume the income schedule is an unrestricted feasible instrument, all consumption is taxed uniformly, there is no evasion or administrative cost, and household preferences depend only on actual consumption and labor, not on the labels attached to payments.

\begin{proposition}[Uniform taxation can be absorbed into the income schedule]
Define
\[
 T'(y)=y-\frac{y-T(y)}{1+t}.
\]
Replacing $(T,t)$ by $(T',0)$ preserves every household's feasible consumption set at every labor income, hence preserves all choices and utility under a common tie rule. At any budget-exhausting choice it also preserves total government receipts.
\end{proposition}
\begin{proof}
Dividing the original budget by $1+t$ gives exactly $p\cdot c\le y-T'(y)$. This equality holds for every $y$, so it preserves the joint labor-consumption opportunity set even without assuming separability. At a budget-exhausting choice, original receipts are
\[
 T(y)+tp\cdot c=T(y)+\frac{t[y-T(y)]}{1+t}
 =\frac{T(y)+ty}{1+t}=T'(y).
\]
Thus public revenue and private utility coincide household by household.
\end{proof}

This is a redundancy result for a uniform tax, not a theorem that differential commodity taxes are always undesirable. It does not establish that a constrained income schedule can implement $T'$, or that hidden labor income and visible consumption are equally enforceable. Legal rate constraints, transfers to nonfilers, labor informality and heterogeneous consumption opportunities can invalidate the replacement. The stronger Atkinson--Stiglitz conclusion about differentiation requires its own information and preference assumptions; it is not inferred from the algebra above.

The result also clarifies a welfare-accounting point. A private compliance burden already subtracted in a household's money-metric utility should not be subtracted again from the government's administrative-resource term. Tax payments and refunds are transfers, while audits, filing time and distorted production can consume real resources. Their treatment depends on the declared public-fund and household weights.

\section{An explicit concealment margin}
Suppose a seller with true taxable sales $s$ chooses concealed sales $h\in[0,s]$. Concealment saves tax $th$, costs real resources $kh^2/2$ with $k>0$, and entails expected monetary penalty $pfh$, where $p\in[0,1]$ is audit probability and $f\ge0$ penalty per hidden sale. The seller maximizes
\[
 (t-pf)h-kh^2/2.
\]
Strict concavity gives the unique choice
\[
 h^*=\min\{s,\max\{0,(t-pf)/k\}\}. \tag{2}
\]
The proof is differentiation on the interior and checking the endpoint derivative signs. Increasing detection can reduce evasion in the interior, but once $pf\ge t$ the concealment margin is zero. If audit cost is $C(p)$, the revenue-maximizing audit probability need not maximize welfare: penalties transfer resources, concealment consumes them, and administrative inputs cost resources.

Invoice matching can alter the effective detection probability or concealment technology. A recorded-sales increase after matching therefore combines behavior and measurement unless true transactions are independently validated. The specific formula (2) is an illustration, not an empirical estimate of a VAT enforcement elasticity.

\section{What missing transactions permit one to infer}
Let $S_{\mathrm{obs}}$ be recorded sales and assume unrecorded taxable sales lie in $[0,M]$. With a uniform rate and no other credits, gross liability lies sharply in $[tS_{\mathrm{obs}},t(S_{\mathrm{obs}}+M)]$: both endpoints are attained by assigning missing sales zero or $M$. Net revenue requires refund, compliance and collection restrictions in addition. If the missing share is allowed to depend on household poverty, observed expenditure shares cannot by themselves identify distributional incidence.

A practical optimization would specify a compact feasible policy set, equilibrium existence and selection, and a validated welfare/revenue map. Uniform approximation errors $|\widehat W-W|\le\epsilon_W$ and $|\widehat R-R|\le\epsilon_R$ support conservative feasibility $\widehat R-C\ge\bar R+\epsilon_R$, but do not ensure that the optimum under this tightened constraint approximates the original optimum without a slack or continuity condition. This guards against treating a noisy revenue constraint as harmless.

The credit-chain identity and uniform-tax replacement are exact within their domains. Joint optimal rates, exemptions, enforcement and redistribution remain unresolved because demand, informality, production and instrument restrictions have not been identified for any actual economy.

\section{Completion Estimate}
40\%

This is a post hoc, heuristic estimate for the full catalogue target.
The attempt proves VAT credit-chain distinctions, exact uniform-tax replacement under unrestricted income schedules, and concealment and missing-sales revenue bounds. Full optimal rate, exemption, and enforcement design still requires identified demand and production responses, legal income-tax limits, invoice-network behavior, and distributional informal-transaction coverage.

\begin{thebibliography}{9}
\bibitem{vat} D. Pomeranz (2015), \emph{No Taxation without Information: Deterrence and Self-Enforcement in the Value Added Tax}, American Economic Review 105(8), 2539--2569. \url{https://www.aeaweb.org/articles?id=10.1257/aer.20130393}. The checked primary record concerns VAT information and enforcement. The accounting and choice models above are self-contained specializations, not a comprehensive optimal-tax theorem attributed to that paper.
\end{thebibliography}

\end{document}